\chapter{Introduction}
\label{sec:intro}

\section{Motivation}

Most self-driving vehicles (SDVs) utilize a 3D object detector to recognize and localize objects in 3D space.
This task is challenging due to occlusion, large intra-class variability, and distant objects, which typically have limited sensor observations.
To overcome these challenges, human drivers rely on their memory.
For example, they may drive more cautiously when remembering a previously observed but now occluded cyclist, who may suddenly enter the road.
\renewcommand{\thefootnote}{\fnsymbol{footnote}}
\footnotetext[2]{Work done while at Waabi}
\renewcommand{\thefootnote}{\arabic{footnote}}

A common approach for improving 3D object detectors is to aggregate a short temporal window of past sensor observations.
Towards this goal, most existing LiDAR-based methods transform a short buffer of sensor data into the current ego (SDV) coordinate frame
to align past and current evidence~\cite{caesar2020nuscenes,sun2022swformer,yang20213d,centerpoint,luo2018fast}. Similarly, camera-based methods stack multiple images ~\cite{zeng2022lift,piergiovanni20214d} as input to existing architectures.
These methods cannot handle long temporal sequences due to computational and memory constraints.
Moreover, temporal stacks of 3D/Bird's-Eye-View (BEV) representations
like point clouds or lifted camera features require a large receptive field,
especially for fast-moving objects~\cite{li2023modar}, further increasing computational burden.

There is a growing interest in long-term temporal fusion.
Scene-level memory approaches~\cite{he2023lef,li2022bevformer,frossard2021strobe}
recurrently fuse scene-level features,
but they can struggle to capture relevant foreground objects.
Other approaches associate objects in memory over time via tracking~\cite{liang2020pnpnet,chen2022mppnet,li2023modar,huang2024ptt}, aggregating past information for each particular object.
However, the associations from the tracker may contain mistakes that can compound over time and lead to information loss. 
Other methods leverage attention from current detection proposals to the past sensor or object information~\cite{he2023msf,yang20213d,hou2024query}. Still, they can be challenging to scale to long histories and suffer from false negatives as the proposals refined into the final detections only come from the present time.

\section{Contributions}

In this paper, we present a simple and sensor-agnostic
add-on for enhancing any existing 3D object detector with long-term memory. We refer to it as \ourmodel{} --- short for \underline{M}emory-\underline{A}ugmented \underline{D}etection, and \cref{fig:hook} illustrates the high-level idea.
\ourmodel{} is a transformer-based model that fuses proposals from a detector with proposals from a memory bank representing past beliefs.
Inspired by recent developments \cite{casas2024detra}, we exploit joint detection and trajectory forecasting.
By storing explicit trajectory forecasts in the memory bank, we can estimate object poses at arbitrary future timestamps for all the objects in the memory.
This enables us to enrich the set of proposals by aligning memory proposals with the current observations.

Training with temporal data can be challenging:
back-propogation through consecutive training examples consumes prohibitive amounts of memory,
training on long sequences can cause over-fitting when back-propagating on every example,
and using memory warm-up can slow down training.
We design a more effective training schedule that begins with short temporal sequences and progressively increases the length, exploiting high data diversity early and closing the gap with inference towards the end.
To ensure the model learns to trust the memory when training on short sequences,
we use cached model outputs from previous training iterations.%we cache model outputs from previous training iterations to make the memory available throughout training.

We demonstrate the generality of our approach by enhancing existing LiDAR-based and camera-based 3D object detection networks with \ourmodel{}, and show considerable improvements over the base detectors on two large-scale datasets: Waymo Open dataset (WOD)~\cite{WOD} and  Argoverse 2 Sensor dataset (AV2)~\cite{wilson2023argoverse}.
Notably, SAFDNet~\cite{safdnet} enhanced with \ourmodel{} achieves state-of-the-art performance on WOD for online detection methods without requiring ensembles or test-time augmentation.

\section{Outline}

This thesis is structured as follows: \cref{sec:related} surveys related work
in 3D object detection and long-term temporal fusion for object detection, 
\cref{sec:method-arch} describes our archiecture and \cref{sec:method-training}
our training procedure, with implementation details in \cref{sec:implementation}.
\cref{sec:experiments} benchmarks our approach against alternatives
\ben{TODO: fill this section out when format is done}